% ============================================================================
% Fichier  : partie2/P02_C05_cobit_principes_architecture.tex
% Titre    : COBIT 2019 : Principes et Architecture
% Auteur   : MINKA MI NGUIDJOÏ Thierry Emmanuel
% Version  : 1.0 -- octobre 2026
% ============================================================================
% Source principale : COBIT 2019 Framework: Introduction and Methodology
% (ISACA, 2018), paraphrasée et traduite par nos soins (dossier Normes).
% ============================================================================

\chapter{COBIT 2019 : Principes et Architecture}
\label{chap:cobit-archi}
\chaptermark{COBIT 2019 : principes et architecture}

\epigraph{%
    « COBIT ne prend aucune décision à votre place. Il dit quelles
    décisions doivent être prises, comment, et par qui. »%
}{D'après ISACA, \textit{COBIT 2019 Framework: Introduction and
Methodology} (2018, notre traduction)}

\niveauChapitre{Fondamental}{%
    Maîtrise du Chapitre~\ref{chap:referentiels} (place de COBIT parmi les
    référentiels) et du Chapitre~\ref{chap:alignement} (trois finalités).%
}

\begin{boxobjectifs}
\begin{itemize}
    \item Retracer l'évolution de COBIT, de la version~1 à COBIT~2019.
    \item Énoncer les six principes du système de gouvernance et les trois
          principes du cadre.
    \item Décrire les sept composantes d'un système de gouvernance.
    \item Présenter les cinq domaines (EDM, APO, BAI, DSS, MEA) et les
          domaines d'intérêt.
\end{itemize}
\end{boxobjectifs}

\bigskip

% ============================================================================
\section*{Positionnement : de la carte au territoire}
% ============================================================================

Le Chapitre~\ref{chap:referentiels} a placé COBIT au centre de la carte des
référentiels~: il fixe les objectifs, et les autres référentiels
fournissent les pratiques. La Partie~\ref{part:cobit}, qui s'ouvre ici,
entre dans ce référentiel central. Elle suit la logique même de COBIT~:
comprendre son architecture (ce chapitre), connaître ses objectifs
(Chapitre~\ref{chap:cobit-objectifs}), concevoir un système de gouvernance
adapté (Chapitre~\ref{chap:cobit-conception}), puis le mesurer et le
mettre en œuvre (Chapitre~\ref{chap:cobit-capacite}).

La question directrice de ce chapitre est simple~: \emph{de quoi un
système de gouvernance des SI est-il fait~?} La réponse de COBIT~2019 tient
en trois mots~: des \textbf{principes}, des \textbf{composantes} et des
\textbf{objectifs}. Le chapitre les présente dans cet ordre, après un
retour sur l'histoire du référentiel.

\separateur

% ============================================================================
\section{Trente Ans d'Évolution}
\label{sec:c05-histoire}
% ============================================================================

\subsection{Du contrôle à la gouvernance}

COBIT est né en 1996 sous le nom de \textit{Control Objectives for
Information and Related Technology}. Il était alors un référentiel
d'objectifs de contrôle à l'usage des auditeurs des systèmes
d'information. Chaque nouvelle version a élargi son public et son champ,
de l'audit vers le management, puis vers la gouvernance de l'ensemble de
l'information et des technologies.

\begin{tableaucours}{L{2.2cm}L{4.4cm}L{5.6cm}}{%
    \tableauHeader{Version} &
    \tableauHeader{Public principal} &
    \tableauHeader{Apport}
}
COBIT 1 (1996) & Auditeurs & Objectifs de contrôle des technologies de
l'information \\
1998, 2000 & Auditeurs, puis managers & Éditions successives~: guides de
management, modèles de maturité \\
COBIT 4.0 et 4.1 (2005, 2007) & Direction informatique, direction
générale & Gouvernance des SI~; 34~processus en quatre domaines (PO, AI,
DS, ME)~\cite{itgi2007cobit41} \\
COBIT 5 (2012) & Ensemble de l'entreprise & Gouvernance de l'entreprise
numérique~; cinq principes, sept « facilitateurs », 37~processus~;
intégration de Val~IT et Risk~IT \\
COBIT 2019 (2018) & Conseil, direction, auditeurs, praticiens & Système
de gouvernance sur mesure~; facteurs de conception, domaines d'intérêt,
40~objectifs, mesure de la capacité alignée sur CMMI \\
\end{tableaucours}
\vspace{-0.8em}
\begin{center}{\small\itshape Tableau~5.1. Évolution de COBIT}\end{center}

\subsection{Ce que COBIT 2019 apporte}

Par rapport à COBIT~5, dont il conserve l'essentiel, COBIT~2019 met en
avant quatre améliorations~\cite{isaca2018cobitintro}~:
\begin{itemize}
    \item \textbf{la flexibilité}~: les facteurs de conception permettent
          d'adapter le système de gouvernance au contexte de
          l'organisation, et l'architecture ouverte permet d'ajouter des
          domaines d'intérêt sans toucher au modèle central~;
    \item \textbf{l'actualité}~: le modèle se rattache plus facilement aux
          normes et réglementations récentes~;
    \item \textbf{l'application prescriptive}~: une fois adapté, le modèle
          sert de prescription pour le système de gouvernance de
          l'organisation~;
    \item \textbf{la mesure de la performance}~: les notions de capacité et
          de maturité sont intégrées au modèle et alignées sur CMMI.
\end{itemize}

\subsection{Ce que COBIT n'est pas}

COBIT~2019 prend soin de dissiper quatre malentendus. COBIT \emph{n'est
pas} une description complète de l'environnement informatique d'une
organisation. Il \emph{n'est pas} un cadre d'organisation des processus
métier. Il \emph{n'est pas} un référentiel technique de gestion des
technologies. Surtout, il \emph{ne prend ni ne prescrit aucune décision}~:
il ne dit pas quelle est la meilleure stratégie, la meilleure architecture
ou le bon niveau de dépense, mais quelles décisions doivent être prises,
comment, et par qui. C'est le sens de l'épigraphe.

\begin{boxalert}
\textbf{Violation typique.} Un cabinet promet à la BANQUE SANAGA de
« mettre la banque en conformité COBIT » en six mois, avec à la clé une
« certification ». La promesse repose sur deux erreurs~: COBIT n'est pas
une norme d'exigences, et il n'existe pas de certification COBIT d'une
organisation (Chapitre~\ref{chap:referentiels}).

\smallskip
\textbf{Règle~:} on ne se « met pas en conformité » avec COBIT~; on s'en
sert pour concevoir et évaluer un système de gouvernance, dont on mesure
la capacité des processus (Chapitre~\ref{chap:cobit-capacite}).
\end{boxalert}

% ============================================================================
\section{Les Principes de COBIT 2019}
\label{sec:c05-principes}
% ============================================================================

COBIT~2019 repose sur deux séries de principes~: six principes qui
décrivent ce qu'un \emph{système de gouvernance} doit être, et trois
principes qui décrivent ce qu'un \emph{cadre de gouvernance}, c'est-à-dire
un référentiel comme COBIT lui-même, doit être~\cite{isaca2018cobitintro}.

\subsection{Six principes pour un système de gouvernance}

\begin{tableaucours}{L{3.6cm}L{8.6cm}}{%
    \tableauHeader{Principe} &
    \tableauHeader{Contenu}
}
1. Créer de la valeur pour les parties prenantes & Toute organisation a
besoin d'un système de gouvernance pour satisfaire ses parties prenantes
et tirer de la valeur des technologies. La valeur est un équilibre entre
bénéfices, risques et ressources. \\
2. Approche globale & Le système de gouvernance est fait de composantes de
types différents, qui fonctionnent ensemble. \\
3. Système dynamique & Chaque fois qu'un facteur de conception change
(stratégie, technologie\ldots), l'effet de ce changement sur le système de
gouvernance doit être examiné. \\
4. Gouvernance distincte du management & Le système distingue clairement
les activités et les structures de gouvernance de celles de management. \\
5. Adapté aux besoins de l'organisation & Le système est ajusté au moyen
de facteurs de conception, qui servent à adapter et à prioriser ses
composantes. \\
6. Couverture de bout en bout & Le système couvre toute l'organisation,
pas seulement la fonction informatique~: tout traitement de l'information
et toute technologie, où qu'ils se trouvent. \\
\end{tableaucours}
\vspace{-0.8em}
\begin{center}{\small\itshape Tableau~5.2. Les six principes d'un système
de gouvernance (d'après COBIT~2019, notre traduction)}\end{center}

\begin{boxscenario}{Opération SANAGA. Les principes violés}
\textbf{Constatation.} L'application « SANAGA Mobile » a été conçue par la
direction marketing avec une agence extérieure, sans passer par la DSI.
Son lancement a fixé la date de bascule du système central.

\medskip
\textbf{Analyse} (confiance élevée). Trois principes au moins ont été
méconnus. Le principe~6~: une application critique s'est développée en
dehors du périmètre de la gouvernance, parce qu'elle ne relevait pas de la
DSI. Le principe~4~: le directeur général a pris seul une décision de
gouvernance (l'arbitrage entre enjeu commercial et risque). Le
principe~3~: le changement de stratégie (le lancement mobile) n'a donné
lieu à aucun réexamen du système de gouvernance du projet HORIZON.
\end{boxscenario}

\subsection{Trois principes pour un cadre de gouvernance}

Les trois autres principes s'appliquent au référentiel lui-même. Un cadre
de gouvernance doit être \textbf{fondé sur un modèle conceptuel}, qui
identifie les composantes et leurs relations, pour garantir la cohérence
et permettre l'automatisation. Il doit être \textbf{ouvert et flexible},
pour accueillir de nouveaux contenus sans perdre sa cohérence. Il doit
être \textbf{aligné sur les grandes normes}, référentiels et
réglementations. Ces principes expliquent pourquoi COBIT renvoie, pour
chaque objectif, aux normes qui le détaillent (Chapitre~\ref{chap:referentiels}).

% ============================================================================
\section{Les Composantes du Système de Gouvernance}
\label{sec:c05-composantes}
% ============================================================================

\subsection{Sept types de composantes}

Pour atteindre ses objectifs, une organisation construit son système de
gouvernance à partir de composantes. COBIT~5 les appelait
« facilitateurs » (\textit{enablers})~; COBIT~2019 parle de
\termecle{composantes}\index{composantes du système de gouvernance}.

\begin{boxdef}
\termecle{Composante}\index{composante} (\textit{component})~: élément qui
contribue, seul et avec les autres, au bon fonctionnement du système de
gouvernance de l'information et des technologies. Les composantes
interagissent et forment un système global. Elles sont de sept types.
\end{boxdef}

\begin{tableaucours}{L{4.2cm}L{8.0cm}}{%
    \tableauHeader{Composante} &
    \tableauHeader{Rôle}
}
Processus & Ensemble organisé de pratiques et d'activités qui atteint un
objectif et produit des résultats. \\
Structures organisationnelles & Instances de décision de l'organisation
(conseil, comités, fonctions). \\
Principes, politiques et procédures & Traduisent le comportement attendu
en consignes pratiques pour le management quotidien. \\
Information & Toute information produite et utilisée~; COBIT s'intéresse à
celle qu'exige le fonctionnement du système de gouvernance. \\
Culture, éthique et comportement & Facteurs souvent sous-estimés de la
réussite de la gouvernance. \\
Personnes, aptitudes et compétences & Nécessaires aux bonnes décisions et
à la bonne exécution. \\
Services, infrastructure et applications & Moyens techniques qui
supportent le système de gouvernance lui-même. \\
\end{tableaucours}
\vspace{-0.8em}
\begin{center}{\small\itshape Tableau~5.3. Les sept composantes (d'après
COBIT~2019, notre traduction)}\end{center}

\begin{center}
\begin{tikzpicture}[
    comp/.style = {
        draw=CouleurBordDef, fill=CouleurDef, rounded corners=4pt,
        text width=2.5cm, minimum height=1.0cm, align=center,
        font=\scriptsize\bfseries\color{CouleurPrimaire}
    },
    centre/.style = {
        draw=CouleurPrimaire, fill=CouleurPrimaire, circle,
        text width=1.9cm, align=center, font=\scriptsize\bfseries\color{white}
    }
]
\node[centre] (c) at (0,0) {Système de\\ gouvernance};
\foreach \ang/\txt in {90/Processus,
    39/{Structures\\ organisationnelles},
    -13/{Principes,\\ politiques,\\ procédures},
    -64/Information,
    -116/{Culture, éthique,\\ comportement},
    -167/{Personnes,\\ compétences},
    141/{Services, infra-\\ structure, applications}}
{
    \node[comp] at (\ang:3.3cm) {\txt};
    \draw[CouleurSecondaire, thick] (c) -- (\ang:2.35cm);
}
\end{tikzpicture}
\end{center}
\vspace{-0.4em}
\begin{center}{\small\itshape Figure~5.1. Les sept composantes du système
de gouvernance}\end{center}

\subsection{La leçon des composantes}

La leçon principale tient au principe~2 (approche globale)~: \textbf{un
processus seul ne suffit jamais}. La BANQUE SANAGA avait, sur le papier,
une procédure de mise en production. Il lui manquait les autres
composantes~: la structure (un comité des changements), l'information (un
procès-verbal de recette fiable), la culture (le droit de dire non à une
date) et les compétences (des testeurs formés au nouveau système).

\begin{boxprocedure}
\textbf{Une composante qui existe vraiment.} Pour chaque composante d'un
objectif prioritaire, l'auditeur cherche une preuve de son existence
\emph{et} de son fonctionnement~:
\begin{itemize}
    \item \emph{processus}~: une procédure approuvée, et des
          enregistrements de son exécution~;
    \item \emph{structure}~: une charte, et des procès-verbaux~;
    \item \emph{politique}~: un document approuvé, daté, diffusé~;
    \item \emph{information}~: le document attendu, à jour et utilisé~;
    \item \emph{culture}~: des signaux observables (alertes remontées,
          décisions de report acceptées)~;
    \item \emph{compétences}~: des fiches de poste et un plan de formation
          suivi~;
    \item \emph{services et applications}~: l'outil qui supporte le
          processus, et ses journaux.
\end{itemize}
\end{boxprocedure}

\subsection{Composantes génériques et variantes}

Les composantes décrites dans le modèle central de COBIT sont
\emph{génériques}~: elles valent en principe pour toute organisation, mais
doivent être adaptées avant d'être mises en œuvre. Une \emph{variante} est
une composante générique adaptée à un usage particulier~: la sécurité de
l'information, DevOps, une réglementation donnée.

% ============================================================================
\section{Objectifs et Domaines}
\label{sec:c05-domaines}
% ============================================================================

\subsection{Objectifs de gouvernance et de management}

\begin{boxdef}
\termecle{Objectif de gouvernance ou de management}\index{objectif de
gouvernance}\index{objectif de management} (\textit{governance or
management objective})~: ce qu'il faut atteindre pour que l'information
et les technologies contribuent aux objectifs de l'organisation. Chaque
objectif se rapporte à un processus de même nom (ou d'un nom proche) et à
un ensemble de composantes des autres types qui aident à l'atteindre.
\end{boxdef}

Un objectif de gouvernance se rapporte à un processus de gouvernance, dont
le conseil et la direction générale répondent. Un objectif de management
se rapporte à un processus de management, qui relève de l'encadrement
supérieur et intermédiaire.

\subsection{Cinq domaines, quarante objectifs}

Les 40~objectifs du modèle central sont regroupés en cinq domaines, dont
les noms sont des verbes~\cite{isaca2018cobitintro}.

\begin{tableaucours}{L{1.4cm}L{3.6cm}cL{5.8cm}}{%
    \tableauHeader{Code} &
    \tableauHeader{Domaine} &
    \tableauHeader{Nb} &
    \tableauHeader{Champ}
}
\multicolumn{4}{l}{\textit{Gouvernance}} \\
EDM & Évaluer, diriger et surveiller & 5 & L'organe de gouvernance évalue
les options stratégiques, oriente la direction sur les options retenues et
surveille la réalisation de la stratégie \\
\multicolumn{4}{l}{\textit{Management}} \\
APO & Aligner, planifier et organiser & 14 & Organisation d'ensemble,
stratégie et activités de soutien \\
BAI & Bâtir, acquérir et implémenter & 11 & Définition, acquisition et mise
en œuvre des solutions, et leur intégration dans les processus métier \\
DSS & Délivrer, servir et soutenir & 6 & Exploitation et soutien des
services, y compris la sécurité \\
MEA & Surveiller, évaluer et apprécier & 4 & Suivi de la performance et de
la conformité aux objectifs internes, au contrôle interne et aux
exigences externes \\
\textbf{Total} & & \textbf{40} & \\
\end{tableaucours}
\vspace{-0.8em}
\begin{center}{\small\itshape Tableau~5.4. Les cinq domaines de COBIT~2019
(intitulés traduits par nos soins)}\end{center}

On retrouve dans ce découpage la distinction des Chapitres~\ref{chap:prologue}
et~\ref{chap:gouv-entreprise}~: le domaine EDM correspond aux trois verbes
de la gouvernance, et les quatre domaines de management au cycle
planifier, bâtir, exploiter, surveiller. Le Chapitre~\ref{chap:cobit-objectifs}
parcourt les 40~objectifs un par un.

\begin{boxobjet}{EDM05 · Assurer l'engagement des parties prenantes}
\textbf{Finalité~:} s'assurer que les parties prenantes sont identifiées et
associées au système de gouvernance, et que la mesure et le compte rendu
de la performance et de la conformité des SI leur sont transparents.

\textbf{Pratiques~:} évaluer les exigences de compte rendu des parties
prenantes~; diriger la communication et le compte rendu~; surveiller la
communication.

\textbf{Sorties~:} principes de compte rendu, plan de communication,
rapports aux parties prenantes.

\textbf{RACI (extrait)~:} \emph{A} au conseil d'administration, \emph{R} au
directeur général, \emph{C} au DSI et au responsable GRC.

\textbf{Indicateurs~:} satisfaction des parties prenantes à l'égard des
comptes rendus~; délai de transmission des rapports exigés par le
régulateur.
\end{boxobjet}
{\small\itshape Intitulé et pratiques résumés et traduits par nos soins à
partir de COBIT~2019~\cite{isaca2018cobitobjectifs}.}

\begin{boxmetier}{Responsable GRC}
\textbf{Les parties prenantes de la BANQUE SANAGA, objectif EDM05}
\begin{enumerate}
    \item Dresser la liste des parties prenantes internes (conseil,
          comités, direction, métiers, audit interne) et externes (COBAC,
          commissaire aux comptes, ANTIC, clients, intégrateurs).
    \item Pour chacune, noter ce qu'elle attend comme compte rendu, à quelle
          fréquence, et sous quelle forme.
    \item Construire le calendrier des comptes rendus de l'année.
    \item Faire approuver ce plan par le comité stratégique des SI.
\end{enumerate}
Après la panne, c'est le rapport à la COBAC qui a révélé l'absence de ce
plan~: personne ne savait qui devait le rédiger, ni avec quelles données.
\end{boxmetier}

% ============================================================================
\section{Domaines d'Intérêt et Facteurs de Conception}
\label{sec:c05-focus}
% ============================================================================

\subsection{Les domaines d'intérêt}

\begin{boxdef}
\termecle{Domaine d'intérêt}\index{domaine d'intérêt} (\textit{focus
area})~: thème, domaine ou question de gouvernance traité par un ensemble
d'objectifs de gouvernance et de management et de leurs composantes.
Exemples donnés par COBIT~2019~: petites et moyennes entreprises,
cybersécurité, transformation numérique, informatique en nuage, protection
de la vie privée, DevOps.
\end{boxdef}

Le nombre de domaines d'intérêt n'est pas limité~: c'est ce qui rend le
modèle ouvert. Un domaine d'intérêt peut combiner des composantes
génériques et des variantes. Pour NOUN PAY, jeune entreprise hébergée dans
le nuage, les domaines d'intérêt « petites et moyennes entreprises » et
« informatique en nuage » sont naturellement pertinents.

\subsection{Les facteurs de conception}

\begin{boxdef}
\termecle{Facteur de conception}\index{facteur de conception}
(\textit{design factor})~: facteur susceptible d'influencer la conception
du système de gouvernance d'une organisation. COBIT~2019 en retient onze~:
la stratégie, les objectifs de l'organisation, le profil de risque, les
problèmes liés aux SI, le paysage des menaces, les exigences de
conformité, le rôle des SI, le mode d'approvisionnement des SI, les
méthodes de mise en œuvre, la stratégie d'adoption des technologies et la
taille de l'organisation.
\end{boxdef}

Les facteurs de conception sont l'outil qui rend COBIT proportionné. Une
banque soumise à la COBAC et une jeune entreprise de 45~salariés n'auront
ni les mêmes objectifs prioritaires, ni les mêmes niveaux de capacité
cibles. Le Chapitre~\ref{chap:cobit-conception} montre comment les
appliquer, et le laboratoire GCI-LAB-02 en fait le calcul.

\begin{boiteRemarque}{Un modèle conçu ailleurs, utilisable ici~?}
Les facteurs de conception répondent en partie à l'objection du
Chapitre~\ref{chap:referentiels}~: les référentiels internationaux
seraient taillés pour d'autres contextes. En faisant de la taille, du
mode d'approvisionnement ou des exigences de conformité des paramètres
explicites, COBIT~2019 permet d'intégrer les réalités camerounaises
(dépendance aux intégrateurs, régulation de la COBAC et de l'ANTIC,
ressources limitées) dans la conception même du système. Encore faut-il
renseigner ces facteurs avec lucidité~: un DSI qui se décrit comme une
organisation « innovante » pour justifier un budget fausse le résultat.
\end{boiteRemarque}

\begin{boxmetier}{Auditeur SI}
\textbf{Vérifier qu'une organisation utilise réellement COBIT}
\begin{enumerate}
    \item Demander le document qui identifie les objectifs prioritaires
          retenus, et la manière dont ils ont été choisis (facteurs de
          conception).
    \item Pour deux ou trois objectifs, vérifier l'existence des sept
          composantes (exigence de contrôle ci-dessus).
    \item Demander la dernière évaluation de capacité et ses suites.
\end{enumerate}
Une organisation qui « applique COBIT » sans pouvoir produire ces trois
pièces en a adopté le vocabulaire, pas la méthode.
\end{boxmetier}

% ============================================================================
\separateur
% ============================================================================

\begin{boxresume}
\textbf{Ce qu'il faut retenir de ce chapitre~:}
\begin{itemize}
    \item \textbf{Du contrôle à la gouvernance}~: né en 1996 pour les
          auditeurs, COBIT est devenu avec la version 2019 un cadre de
          gouvernance et de management sur mesure.
    \item \textbf{COBIT ne décide pas}~: il dit quelles décisions doivent
          être prises, comment et par qui. Il n'est ni une norme
          certifiable, ni un référentiel technique.
    \item \textbf{Six principes pour le système}~: valeur pour les parties
          prenantes, approche globale, système dynamique, gouvernance
          distincte du management, adaptation, couverture de bout en bout.
    \item \textbf{Trois principes pour le cadre}~: modèle conceptuel,
          ouverture, alignement sur les normes.
    \item \textbf{Sept composantes}~: processus, structures, politiques,
          information, culture, compétences, services et applications. Un
          processus seul ne suffit jamais.
    \item \textbf{Cinq domaines, quarante objectifs}~: EDM (5) pour la
          gouvernance~; APO (14), BAI (11), DSS (6) et MEA (4) pour le
          management.
    \item \textbf{Domaines d'intérêt et facteurs de conception} rendent le
          modèle ouvert et proportionné.
\end{itemize}
\end{boxresume}

\bigskip

\begin{boxexo}[title={\ding{45}\quad Exercice~5.1. Classer par domaine \niveaubase}]
Rattachez chacune des activités suivantes au domaine COBIT qui convient
(EDM, APO, BAI, DSS, MEA), et justifiez en une phrase.
\begin{enumerate}[(a)]
    \item Le conseil fixe le niveau de risque informatique acceptable.
    \item La DSI élabore le plan stratégique des SI à trois ans.
    \item Une équipe recette le nouveau logiciel de crédit.
    \item Le centre de services traite une panne de messagerie.
    \item L'audit interne évalue les contrôles d'accès.
    \item La direction des achats négocie le contrat d'un infogérant.
\end{enumerate}
\end{boxexo}

\begin{boxexo}[title={\ding{45}\quad Exercice~5.2. Les composantes chez KRIBILOG \niveaumoyen}]
M.~MBARGA, responsable des services informatiques de KRIBILOG SA, veut
améliorer la gestion des incidents du logiciel de suivi des conteneurs.
\begin{enumerate}[(a)]
    \item Pour chacune des sept composantes, décrivez ce qui devrait
          exister chez KRIBILOG.
    \item Pour chacune, indiquez la pièce qui prouverait son existence et
          son fonctionnement.
    \item Lesquelles relèvent de KRIBILOG, et lesquelles de l'infogérant
          NETPORT SERVICES~? Qu'en déduisez-vous pour le contrat~?
\end{enumerate}
\end{boxexo}

\begin{boxexo}[title={\ding{45}\quad Exercice~5.3. De COBIT 4.1 à COBIT 2019 à l'ANPP \niveauexpert}]
Les procédures informatiques de l'ANPP (opération MEFOU) ont été rédigées
en 2010 à partir de COBIT~4.1 et de ses quatre domaines (PO, AI, DS, ME).
\begin{enumerate}[(a)]
    \item Établissez la correspondance entre les quatre domaines de
          COBIT~4.1 et les cinq domaines de COBIT~2019. Lequel n'a pas
          d'équivalent direct, et pourquoi est-ce important pour l'ANPP~?
    \item À l'aide des documents du dossier \texttt{Normes}, identifiez
          deux objectifs de COBIT~2019 sans équivalent clair dans
          COBIT~4.1, et expliquez leur intérêt pour l'agence.
    \item Rédigez une note de deux pages au directeur général proposant
          une démarche de mise à jour, en appliquant les principes~3 et~5.
\end{enumerate}
\end{boxexo}

\bigskip

\begin{boxinfo}
\textbf{Transition.} Ce chapitre a présenté l'architecture de COBIT~2019~:
principes, composantes, domaines. Le Chapitre~\ref{chap:cobit-objectifs}
entre dans le détail des 40~objectifs, apprend à lire leur structure et
montre comment ils s'enchaînent depuis les objectifs de l'organisation
par la cascade des objectifs.
\end{boxinfo}
